\documentclass[10pt,a4paper]{amsart}
\usepackage{amsmath,amssymb,amsthm,graphicx}
% Retained preamble for the Stacks Project chapter morphisms.tex (Morphisms of Schemes).
% The chapter file pinned-original-source.tex carries only \input{preamble} on its first line, so no
% preamble text lives inside the chapter itself.  The Stacks Project uses one single shared preamble for
% every chapter; the snapshot stages it verbatim in sources/stacks-chow/preamble.tex.  The block below is
% copied mechanically, line for line, from that file's lines 14-18 (the Xy-pic setup) and lines 154-217
% (the native theorem declarations together with the mathematical macros).  Nothing was added, removed or
% redefined.  The wrapper class and the amsmath/amssymb/amsthm/graphicx base packages are supplied by the
% assembler, so the class/\externaldocument part of the shared preamble is not repeated here.
% For commutative diagrams we use Xy-pic
\usepackage[all]{xy}

% We use 2cell for 2-commutative diagrams.
\xyoption{2cell}
% Theorem environments.
%
\theoremstyle{plain}
\newtheorem{theorem}[subsection]{Theorem}
\newtheorem{proposition}[subsection]{Proposition}
\newtheorem{lemma}[subsection]{Lemma}

\theoremstyle{definition}
\newtheorem{definition}[subsection]{Definition}
\newtheorem{example}[subsection]{Example}
\newtheorem{exercise}[subsection]{Exercise}
\newtheorem{situation}[subsection]{Situation}

\theoremstyle{remark}
\newtheorem{remark}[subsection]{Remark}
\newtheorem{remarks}[subsection]{Remarks}

\numberwithin{equation}{subsection}

% Macros
%
\def\lim{\mathop{\mathrm{lim}}\nolimits}
\def\colim{\mathop{\mathrm{colim}}\nolimits}
\def\Spec{\mathop{\mathrm{Spec}}}
\def\Hom{\mathop{\mathrm{Hom}}\nolimits}
\def\Ext{\mathop{\mathrm{Ext}}\nolimits}
\def\SheafHom{\mathop{\mathcal{H}\!\mathit{om}}\nolimits}
\def\SheafExt{\mathop{\mathcal{E}\!\mathit{xt}}\nolimits}
\def\Sch{\mathit{Sch}}
\def\Mor{\mathop{\mathrm{Mor}}\nolimits}
\def\Ob{\mathop{\mathrm{Ob}}\nolimits}
\def\Sh{\mathop{\mathit{Sh}}\nolimits}
\def\NL{\mathop{N\!L}\nolimits}
\def\CH{\mathop{\mathrm{CH}}\nolimits}
\def\proetale{{pro\text{-}\acute{e}tale}}
\def\etale{{\acute{e}tale}}
\def\QCoh{\mathit{QCoh}}
\def\Ker{\mathop{\mathrm{Ker}}}
\def\Im{\mathop{\mathrm{Im}}}
\def\Coker{\mathop{\mathrm{Coker}}}
\def\Coim{\mathop{\mathrm{Coim}}}

% Boxtimes
%
\DeclareMathSymbol{\boxtimes}{\mathbin}{AMSa}{"02}

%
% Macros for moduli stacks/spaces
%
\def\QCohstack{\mathcal{QC}\!\mathit{oh}}
\def\Cohstack{\mathcal{C}\!\mathit{oh}}
\def\Spacesstack{\mathcal{S}\!\mathit{paces}}
\def\Quotfunctor{\mathrm{Quot}}
\def\Hilbfunctor{\mathrm{Hilb}}
\def\Curvesstack{\mathcal{C}\!\mathit{urves}}
\def\Polarizedstack{\mathcal{P}\!\mathit{olarized}}
\def\Complexesstack{\mathcal{C}\!\mathit{omplexes}}
% \Pic is the operator that assigns to X its picard group, usage \Pic(X)
% \Picardstack_{X/B} denotes the Picard stack of X over B
% \Picardfunctor_{X/B} denotes the Picard functor of X over B
\def\Pic{\mathop{\mathrm{Pic}}\nolimits}
\def\Picardstack{\mathcal{P}\!\mathit{ic}}
\def\Picardfunctor{\mathrm{Pic}}
\def\Deformationcategory{\mathcal{D}\!\mathit{ef}}

% --- appended 2026-10-12 (Chief): the project own definitions of the reference/slogan/history
% display environments, copied mechanically from the shared Stacks preamble lines 8-12
% (\usepackage{verbatim} plus the three \newenvironment lines, which the chapter preamble keeps
% verbatim from the project preamble). Nothing added, removed or redefined: these environments
% carry no mathematical content and the project itself typesets them as comments. Without these
% three lines a counted statement whose own text uses them fails with Environment slogan undefined.
\usepackage{verbatim}
\newenvironment{reference}{\comment}{\endcomment}
\newenvironment{slogan}{\comment}{\endcomment}
\newenvironment{history}{\comment}{\endcomment}
\usepackage{cleveref}
\crefname{originallemma}{original lemma}{original lemmas}
\Crefname{originallemma}{Original lemma}{Original lemmas}
\makeatletter
\newlabel{algebra-lemma-characterize-finite-in-terms-of-integral}{{02JJ}{source}{}{sourcepointer-1}{}}
\newlabel{algebra-lemma-characterize-finite-in-terms-of-integral@cref}{{[originallemma][0][]02JJ}{[1][0][]source}{}{}{}}
\makeatother
\crefname{originallemma}{original lemma}{original lemmas}
\Crefname{originallemma}{Original lemma}{Original lemmas}
\makeatletter
\newlabel{algebra-lemma-support-closed}{{00L2}{source}{}{sourcepointer-2}{}}
\newlabel{algebra-lemma-support-closed@cref}{{[originallemma][0][]00L2}{[1][0][]source}{}{}{}}
\makeatother
\crefname{originallemma}{original lemma}{original lemmas}
\Crefname{originallemma}{Original lemma}{Original lemmas}
\makeatletter
\newlabel{algebra-lemma-support-zero}{{0585}{source}{}{sourcepointer-3}{}}
\newlabel{algebra-lemma-support-zero@cref}{{[originallemma][0][]0585}{[1][0][]source}{}{}{}}
\makeatother
\crefname{originallemma}{original lemma}{original lemmas}
\Crefname{originallemma}{Original lemma}{Original lemmas}
\makeatletter
\newlabel{lemma-integral-universally-closed}{{01WM}{source}{}{sourcepointer-4}{}}
\newlabel{lemma-integral-universally-closed@cref}{{[originallemma][0][]01WM}{[1][0][]source}{}{}{}}
\makeatother
\crefname{originallemma}{original lemma}{original lemmas}
\Crefname{originallemma}{Original lemma}{Original lemmas}
\makeatletter
\newlabel{lemma-nth-and-nplusone-implies-square-and-cube}{{0CNA}{source}{}{sourcepointer-5}{}}
\newlabel{lemma-nth-and-nplusone-implies-square-and-cube@cref}{{[originallemma][0][]0CNA}{[1][0][]source}{}{}{}}
\makeatother
\crefname{originallemma}{original lemma}{original lemmas}
\Crefname{originallemma}{Original lemma}{Original lemmas}
\makeatletter
\newlabel{lemma-special-elements-and-localization}{{0CN9}{source}{}{sourcepointer-6}{}}
\newlabel{lemma-special-elements-and-localization@cref}{{[originallemma][0][]0CN9}{[1][0][]source}{}{}{}}
\makeatother
\crefname{originallemma}{original lemma}{original lemmas}
\Crefname{originallemma}{Original lemma}{Original lemmas}
\makeatletter
\newlabel{lemma-subalgebra-inherits}{{0CN7}{source}{}{sourcepointer-7}{}}
\newlabel{lemma-subalgebra-inherits@cref}{{[originallemma][0][]0CN7}{[1][0][]source}{}{}{}}
\makeatother
\title{An extension which is a universal homeomorphism on spectra onto isomorphic residue fields contains an element outside the base whose square and cube lie in the base}
\author{The Stacks Project Authors}
\date{stacks-project:morphisms.tex:local-snapshot-20261012; GNU FDL 1.2 or later}
\begin{document}
\maketitle
\noindent\textit{Editorial note.} All delivered mathematical text is the author's own, copied from the pinned Stacks Project chapter file pinned-original-source.tex; no statement or proof step was rewritten, reordered or removed.\par\medskip
\noindent\textit{Editorial note.} Following the multi-claim convention for this source, source\_url is the source-level URL https://stacks.math.columbia.edu/ and the source label lemma-square-and-cube is carried in claim\_key and here rather than in a per-claim tag URL.\par\medskip
\noindent\textit{Editorial note.} The preamble is the chapter's shared Stacks Project preamble, retained as sources/stacks-morphisms/retained-preamble.tex: a line-for-line excerpt of the shared Stacks preamble staged with this snapshot. No macro was invented or redefined, and this card adds no preamble text of its own.\par\medskip
\noindent\textit{Editorial note.} Statement span 11614-11620 and proof span 11622-11654 come from the halfA pool card operations/pool-cards/stacks-morphisms-halfA.json and were re-read in the pinned chapter file before this card was written.\par\medskip
\noindent\textit{Editorial note.} Tag self-check per LANE\_BRIEF 10.5 (5): the label lemma-square-and-cube resolves in the project's own tag list to canonical tag 0CNB (line 15886 of the index, full label morphisms-lemma-square-and-cube). reference/public-items.jsonl holds no row whose primary source\_file is morphisms.tex, so the intersection with the same-chapter public tag ids is empty and the candidate passes.\par\medskip
\noindent\textit{Editorial note.} The gate was run as: python operations/agent-authoring/tools/precard\_gate.py --source morphisms.tex --label lemma-square-and-cube --start 11614 --end 11620 --proof-start 11622 --proof-end 11654 --local-id 34094, giving RESULT=OK with all four span guards passing, no public interval row for this chapter and occupancy FREE for morphisms.tex 11614 11654. Spans were recomputed live from the label in the pinned chapter file and the located text matched the pool byte counts. The interval lies outside the two windows of this chapter handed to lane-s1 (8000-10000 and 12000-14000).\par\medskip
\noindent\textit{Editorial note.} 7 cross-references occur inside the counted spans; they are listed in external\_references as navigation pointers only.\par\medskip
\noindent\textit{Allowed standard prerequisites (editorial).}\begin{itemize}
\item A ring extension inducing a universal homeomorphism on spectra is integral, cited from the same chapter. Cited at proof lines 11623-11624. Gives that the extension is integral, which is what allows the reduction to a single generator and finiteness.
\item Passing to a subalgebra generated by one element when testing the conclusion, cited from the same chapter. Cited at proof lines 11625-11627. Reduces the search for an element outside the base to the case where the extension is generated by one element.
\item An integral algebra generated by one element is finite, and the support of a finite module is closed and nonempty when the module is nonzero, cited from the algebra chapter. Cited at proof lines 11627-11631. Makes B/A a finite module with nonempty closed support, so that it makes sense to localize at a minimal prime of that support.
\item The conclusion is preserved by localizing the extension at a prime, cited from the same chapter. Cited at proof lines 11632-11635. Allows the reduction to a local base ring with B finite and the support of B/A equal to the maximal ideal.
\item Two consecutive powers of an element lying in the base imply that its square and cube lie in the base, cited from the same chapter. Cited at proof lines 11647-11649. Converts the statement that f to the n and f to the n+1 both lie in A into the conclusion of the lemma.
\end{itemize}\medskip
\ifdefined\hypertarget\hypertarget{sourcepointer-1}{}\fi
\noindent\textit{Source reference.} Original lemma 02JJ: focused-seeds/AGP/sources/stacks-snapshot-20261011/tags.txt, source lines 3338--3338. Same-chapter reference used in the retained author text at source line(s) 11029, 11628, 11676. It is supplied only as a navigation pointer so that the retained text compiles; the counted item delivers the author's own statement and the author's complete proof, and no part of the referenced item is copied in or substituted for the counted argument.\par\medskip
\ifdefined\hypertarget\hypertarget{sourcepointer-2}{}\fi
\noindent\textit{Source reference.} Original lemma 00L2: focused-seeds/AGP/sources/stacks-snapshot-20261011/tags.txt, source lines 772--772. Same-chapter reference used in the retained author text at source line(s) 705, 759, 11631, 11640, 11679, 11688. It is supplied only as a navigation pointer so that the retained text compiles; the counted item delivers the author's own statement and the author's complete proof, and no part of the referenced item is copied in or substituted for the counted argument.\par\medskip
\ifdefined\hypertarget\hypertarget{sourcepointer-3}{}\fi
\noindent\textit{Source reference.} Original lemma 0585: focused-seeds/AGP/sources/stacks-snapshot-20261011/tags.txt, source lines 6768--6768. Same-chapter reference used in the retained author text at source line(s) 11631, 11679. It is supplied only as a navigation pointer so that the retained text compiles; the counted item delivers the author's own statement and the author's complete proof, and no part of the referenced item is copied in or substituted for the counted argument.\par\medskip
\ifdefined\hypertarget\hypertarget{sourcepointer-4}{}\fi
\noindent\textit{Source reference.} Original lemma 01WM: focused-seeds/AGP/sources/stacks-snapshot-20261011/tags.txt, source lines 2487--2487. Same-chapter reference used in the retained author text at source line(s) 10976, 11172, 11375, 11380, 11512, 11552, 11624, 11672, 11769, 14152, 14165, 14544. It is supplied only as a navigation pointer so that the retained text compiles; the counted item delivers the author's own statement and the author's complete proof, and no part of the referenced item is copied in or substituted for the counted argument.\par\medskip
\ifdefined\hypertarget\hypertarget{sourcepointer-5}{}\fi
\noindent\textit{Source reference.} Original lemma 0CNA: focused-seeds/AGP/sources/stacks-snapshot-20261011/tags.txt, source lines 15885--15885. Same-chapter reference used in the retained author text at source line(s) 11648, 11696. It is supplied only as a navigation pointer so that the retained text compiles; the counted item delivers the author's own statement and the author's complete proof, and no part of the referenced item is copied in or substituted for the counted argument.\par\medskip
\ifdefined\hypertarget\hypertarget{sourcepointer-6}{}\fi
\noindent\textit{Source reference.} Original lemma 0CN9: focused-seeds/AGP/sources/stacks-snapshot-20261011/tags.txt, source lines 15884--15884. Same-chapter reference used in the retained author text at source line(s) 11635, 11683, 14852. It is supplied only as a navigation pointer so that the retained text compiles; the counted item delivers the author's own statement and the author's complete proof, and no part of the referenced item is copied in or substituted for the counted argument.\par\medskip
\ifdefined\hypertarget\hypertarget{sourcepointer-7}{}\fi
\noindent\textit{Source reference.} Original lemma 0CN7: focused-seeds/AGP/sources/stacks-snapshot-20261011/tags.txt, source lines 15882--15882. Same-chapter reference used in the retained author text at source line(s) 11625, 11673, 11931. It is supplied only as a navigation pointer so that the retained text compiles; the counted item delivers the author's own statement and the author's complete proof, and no part of the referenced item is copied in or substituted for the counted argument.\par\medskip
\begin{lemma}
\label{lemma-square-and-cube}
Let $A \subset B$ be a ring extension such that $\Spec(B) \to \Spec(A)$
is a universal homeomorphism inducing isomorphisms on residue fields.
If $A \not = B$, then there exists a $b \in B$, $b \not \in A$ with
$b^2 \in A$ and $b^3 \in A$.
\end{lemma}

\begin{proof}
Recall that $A \subset B$ is integral
(Lemma \ref{lemma-integral-universally-closed}).
By Lemma \ref{lemma-subalgebra-inherits}
we may assume that $B$ is generated by a single element
over $A$. Hence $B$ is finite over $A$
(Algebra, Lemma \ref{algebra-lemma-characterize-finite-in-terms-of-integral}).
Hence the support of $B/A$ as an $A$-module is
closed and not empty (Algebra, Lemmas
\ref{algebra-lemma-support-closed} and \ref{algebra-lemma-support-zero}).
Let $\mathfrak p \subset A$ be a minimal prime
of the support. After replacing $A \subset B$ by
$A_\mathfrak p \subset B_\mathfrak p$ (permissible by
Lemma \ref{lemma-special-elements-and-localization})
we may assume that $(A, \mathfrak m)$ is a local ring,
that $B$ is finite over $A$, and that $B/A$ has support $\{\mathfrak m\}$
as an $A$-module. Since $B/A$ is a finite module,
we see that $I = \text{Ann}_A(B/A)$ satisfies $\mathfrak m = \sqrt{I}$
(Algebra, Lemma \ref{algebra-lemma-support-closed}).
Let $\mathfrak m' \subset B$ be the unique prime ideal lying over
$\mathfrak m$. Because $\Spec(B) \to \Spec(A)$ is a homeomorphism,
we find that $\mathfrak m' = \sqrt{IB}$.
For $f \in \mathfrak m'$ pick $n \geq 1$ such that
$f^n \in IB$. Then also $f^{n + 1} \in IB$.
Since $IB \subset A$ by our choice of $I$ we conclude that
$f^n, f^{n + 1} \in A$. Using
Lemma \ref{lemma-nth-and-nplusone-implies-square-and-cube}
we conclude our lemma is true if $\mathfrak m' \not \subset A$.
However, if $\mathfrak m' \subset A$, then $\mathfrak m' = \mathfrak m$
and we conclude that $A = B$ as the residue fields
are isomorphic as well by assumption. This contradiction
finishes the proof.
\end{proof}

\end{document}
